\chapter{Conclusions}
\label{ch:conclusions}

The power stage of a three-phase two-level SiC inverter for \SI{800}{\volt},
\SI{30}{\Arms} and \SI{20}{\kilo\hertz} has been designed, with its cooling
and the schematic and layout of the module. This chapter sets the results
against the brief, names what remains open, and says where the design would
go next.

\section{Results against the brief}

\begin{table}[htbp]
\centering
\caption{The design against the brief and against the device criteria of
\cref{tab:criteria}.}
\label{tab:results}
\small
\begin{tabularx}{\textwidth}{L{0.26\textwidth}YL{0.2\textwidth}}
\toprule
Asked & Achieved & Where \\
\midrule
\SI{800}{\volt}, \SI{30}{\Arms} at any power factor &
  \SI{25.48}{\kilo\watt}, \SI{128.1}{\watt} of loss, efficiency \SI{99.50}{\percent} at the rated point;
  \SIrange{127.5}{128.1}{\watt} for $\cos\varphi \leq 0.87$ and \SI{131.5}{\watt} at unity & \cref{ch:losses}, \cref{sec:pfsweep} \\
\SI{20}{\kilo\hertz} &
  \SIrange{90.5}{166.1}{\watt} and $\Tj$ \SIrange{76.9}{86.8}{\celsius} over \SIrange{10}{30}{\kilo\hertz} &
  \cref{ch:losses} \\
Negligible output current ripple &
  \SI{2.46}{\ampere} peak-to-peak, \SI{0.84}{\percent} in RMS terms, with \SI{2.5}{\milli\henry} of load &
  \cref{ch:modulation} \\
DC-link voltage ripple $< \SI{0.5}{\percent}$, worst case &
  \SI{0.17}{\percent} at the rated point; \SI{2.86}{\volt}, \SI{0.36}{\percent}, in the worst case with \SI{200}{\micro\farad} &
  \cref{sec:dclink}, \cref{ch:model} \\
Forced air or liquid, \SI{50}{\celsius} ambient &
  liquid, six-port rail and AlN pad: \SI{77.1}{\celsius} at the rated point, \SI{97.9}{\celsius} at locked rotor, \SI{117.8}{\celsius} with all tolerances, against \SI{150}{\celsius};
  one fan with the same pad: \SI{108.3}{} and \SI{129.1}{\celsius} &
  \cref{sec:cooling} \\
\midrule
C1, locked rotor $\leq \SI{120}{\celsius}$ & \SI{97.9}{\celsius} & \cref{sec:coolalt} \\
C2, die $\leq \SI{1080}{\volt}$ at \SI{950}{\volt}, \SI{60}{\ampere} &
  \SI{1073.5}{\volt} at the \SI{20}{\nano\henry} budget; \SI{1040}{\volt} with the loop extracted from the copper &
  \cref{sec:kelvin}, \cref{sec:loopextract} \\
C3, drain--source creepage at \SI{900}{\volt} &
  \SI{3.24}{\milli\metre} against \SI{2.80}{\milli\metre} in PD1, with conformal coating &
  \cref{sec:insulation} \\
C4, automotive & power device AEC-Q101 qualified & \cref{sec:deviceselection} \\
C5, declared $t_\mathrm{SC} \geq \SI{1.5}{\micro\second}$ & \textbf{open}: not declared by the manufacturer & \cref{app:assumptions}, A.3 \\
\bottomrule
\end{tabularx}
\end{table}

The brief is met wherever the design can show it, with one open point: the
short-circuit withstand time, which the manufacturer does not declare and
which no simulation can replace. The stalled motor, a requirement this
design added, sizes the design. It decided the device and the pad, and it is
met with \SI{22}{\celsius} of margin on the criterion once the heat path is
designed rather than assumed.

\section{What remains open}

Three assumptions carry the thermal and the switching results, and each is
listed in \cref{app:assumptions} with the way to close it. The case-to-sink
resistance of the mounted pad (A.5) is the largest single uncertainty. The
loop inductance (A.7) is extracted from the copper, not measured, and the
gate resistance was chosen at \SI{20}{\nano\henry}, \SI{12}{\nano\henry}
above the extracted value. The switching energies (A.4) come from the
datasheet because the model underestimates them by a factor of 3.4 at turn-on.

\section{Possible developments}

\paragraph{Discontinuous modulation at standstill.} Clamping the phase that
carries the peak current to the positive rail lowers the hottest junction at
locked rotor by about \SI{12}{\celsius} on the real rail with the adopted pad
and the converter loss by a third (\cref{sec:lrdpwm}). It stays a
recommendation to the controller (\cref{tab:sysassume}), because hardware is
not sized on a control strategy the module cannot guarantee.

\paragraph{Internally isolated packages.} A TO-247-4 with the ceramic inside,
under the die, removes the external pad and one interface with the same
footprint (\cref{sec:coolalt}). When one is offered at this resistance with
automotive qualification it is the natural next step for the power devices.

\paragraph{A power module.} The discrete TO-247-4L sets two ceilings. The
first is inductance: \SI{14.6}{\nano\henry} of package contribution that no
layout can remove (\cref{sec:kelvin}), which forces
$R_{g,\mathrm{off}}$ to \SI{12.5}{\ohm} and makes switching more than half
of the loss budget. The second is insulation: the \SI{5.08}{\milli\metre}
lead pitch leaves \SI{3.24}{\milli\metre} of creepage between drain and
source, which passes with \SI{16}{\percent} of margin only under conformal
coating, and which IEC 60664-4~\cite{iec606644} would derate above
\SI{30}{\kilo\hertz}, the upper end of the sweep of \cref{tab:sweep-fsw}. A
power module with an integrated busbar lifts both ceilings together. The
selection study of \cref{sec:deviceselection} found none that fits at
\SI{30}{\Arms}: the small six-packs are industrial and run hot, and the
automotive single-switch modules are meant to be sintered onto a cooler. At a
few times the current a module would be the only answer.
